\section{Locked Final-Week Evaluation}
\label{sec:final}

% Full-width tables appear at the earliest on the page after the one where they are
% defined, so they are defined here, ahead of the text that discusses them.
\begin{table*}[t]
\centering\footnotesize
\caption{Mean Newsvendor Cost of Every Locked Policy for A-Trained Models on the A History (Lowest Cost per Panel and Setting in Bold)}
\label{tab:policies}
\begin{tabular}{@{}llrrrr@{}}
\toprule
 & & & \multicolumn{3}{c}{Mean cost} \\
\cmidrule(l){4-6}
Policy & Population & $n$ & $\beta=0.25$ & $\beta=0.50$ & $\beta=0.75$ \\
\midrule
\multicolumn{6}{@{}l}{\emph{(a) Full-sample exact-donor candidates}} \\
Seasonal naive + residual quantile & Full sample & 20,641 & 0.6759 & 0.4430 & 0.7853 \\
Raw-sales LightGBM + residual quantile & Full sample & 20,641 & \textbf{0.5448} & \textbf{0.3868} & 0.7366 \\
Raw-sales quantile LightGBM & Full sample & 20,641 & 0.5749 & 0.3999 & 0.7503 \\
Demand-quantile LightGBM (uncalibrated) & Full sample & 20,641 & 0.5816 & 0.3946 & 0.7152 \\
Demand-quantile LightGBM + pooled conformal & Full sample & 20,641 & 0.5766 & 0.3891 & \textbf{0.6831} \\
Kaplan--Meier order percentile & Full sample & 20,641 & 0.8404 & 0.6520 & 1.3878 \\
Demand-quantile + adaptive calibration (sensitivity) & Full sample & 20,641 & 0.5815 & 0.3932 & 0.6963 \\
\midrule
\multicolumn{6}{@{}l}{\emph{(b) Declared deep subset of 600 series (separate population)}} \\
TFT (direct) + residual quantile & Deep subset & 2,510 & 0.6892 & 0.4949 & 0.9483 \\
TimesNet$\rightarrow$TFT + residual quantile & Deep subset & 2,510 & \textbf{0.6586} & \textbf{0.4687} & \textbf{0.8759} \\
\bottomrule
\end{tabular}

\end{table*}

\begin{table*}[t]
\centering\footnotesize
\caption{Pre-Specified Paired Contrasts for the Pooled-Conformal Policy over the Three Matched Histories (Negative Values Favor Pooled Conformal)}
\label{tab:contrasts}
\begin{tabular}{@{}lllrrr@{}}
\toprule
Pooled conformal minus & Cost setting & Population & $n$ & Difference & 95\% percentile interval \\
\midrule
Uncalibrated demand quantile & Excess stock more costly & Matched full sample & 61,923 & $-0.0052$ & $[-0.0068, -0.0034]$ \\
Uncalibrated demand quantile & Equal costs & Matched full sample & 61,923 & $-0.0077$ & $[-0.0096, -0.0062]$ \\
Uncalibrated demand quantile & Shortage more costly & Matched full sample & 61,923 & $-0.0420$ & $[-0.0483, -0.0365]$ \\
\addlinespace
Raw-sales LightGBM + residual & Excess stock more costly & Matched full sample & 61,923 & $+0.0208$ & $[0.0016, 0.0432]$ \\
Raw-sales LightGBM + residual & Equal costs & Matched full sample & 61,923 & $-0.0091$ & $[-0.0204, 0.0052]$ \\
Raw-sales LightGBM + residual & Shortage more costly & Matched full sample & 61,923 & $-0.0790$ & $[-0.0946, -0.0615]$ \\
\addlinespace
Seasonal naive + residual & Excess stock more costly & Matched full sample & 61,923 & $-0.1213$ & $[-0.1465, -0.0876]$ \\
Seasonal naive + residual & Equal costs & Matched full sample & 61,923 & $-0.0694$ & $[-0.0872, -0.0439]$ \\
Seasonal naive + residual & Shortage more costly & Matched full sample & 61,923 & $-0.1376$ & $[-0.1723, -0.0902]$ \\
\addlinespace
Kaplan--Meier percentile & Excess stock more costly & KM-supported rows & 61,523 & $-0.4043$ & $[-0.4276, -0.3824]$ \\
Kaplan--Meier percentile & Equal costs & KM-supported rows & 58,183 & $-0.3223$ & $[-0.3429, -0.3046]$ \\
Kaplan--Meier percentile & Shortage more costly & KM-supported rows & 47,294 & $-0.7475$ & $[-0.8175, -0.6812]$ \\
\bottomrule
\end{tabular}

\end{table*}

\begin{table*}[t]
\centering\footnotesize
\caption{Transfer Gaps of the Pooled-Conformal Policy When Shortage Is More Costly}
\label{tab:transfer}
\begin{tabular}{@{}cclrrrrr@{}}
\toprule
Train & Test history & Population & $n$ & Matched & Unseen & Gap & 95\% percentile interval \\
\midrule
A & B & Full sample & 20,641 & 0.6831 & 0.6932 & $+0.0101$ & $[0.0076, 0.0128]$ \\
A & C & Full sample & 20,641 & 0.6831 & 0.7303 & $+0.0472$ & $[0.0391, 0.0554]$ \\
B & A & Full sample & 20,641 & 0.6940 & 0.6908 & $-0.0032$ & $[-0.0059, -0.0006]$ \\
B & C & Full sample & 20,641 & 0.6940 & 0.7306 & $+0.0365$ & $[0.0275, 0.0447]$ \\
C & A & Full sample & 20,641 & 0.7248 & 0.7120 & $-0.0128$ & $[-0.0219, -0.0050]$ \\
C & B & Full sample & 20,641 & 0.7248 & 0.7086 & $-0.0162$ & $[-0.0242, -0.0083]$ \\
\bottomrule
\end{tabular}

\end{table*}

The final evaluation scores donor rows only. Each of the nine transfer cells, one per pair of training mechanism and evaluation history, holds the same 20,641 final donor series-dates, on which the six full-sample policies other than Kaplan--Meier are scored. Kaplan--Meier scores only rows with an estimable order, between 6,012 and 20,641 per cell, and the two neural policies score the 2,510 final donor rows of the deep subset; no predictions are imputed for the remaining 4,400 series.

\subsection{Policy comparison}
\label{sec:policies}

Table~\ref{tab:policies} gives the mean costs of the A-trained models on the A history. The best full-sample policy depends on the cost ratio: the raw-sales LightGBM with residual quantiles is cheapest when excess stock is more costly and, narrowly, at equal costs (0.3868 against 0.3891 for pooled conformal), whereas the pooled-conformal policy is cheapest when shortage is more costly. The B history gives the same ranking; on the C history the pooled-conformal policy is already the cheaper of the two at equal costs (0.4094 against 0.4411). Kaplan--Meier is the most expensive full-sample policy in every setting, and within the deep subset TFT on the recovered history costs less than direct TFT in every matched history and setting.

Table~\ref{tab:contrasts} gives the paired contrasts that the lock pre-specified, pooled over the three matched histories. Calibration lowers cost relative to the uncalibrated demand quantiles in all three settings, and the pooled-conformal policy is cheaper than the seasonal-naive and Kaplan--Meier policies throughout. Against the raw-sales LightGBM with residual quantiles the sign depends on the cost ratio: pooled conformal costs more when excess stock is more costly, the equal-cost interval contains zero, and it costs less when shortage is more costly. Moving from $\beta=0.25$ to $\beta=0.75$ raises the pooled-conformal fill rate on the A history from 0.614 to 0.891 and the leftover share from 0.056 to 0.226, and the C history is the most expensive in every setting. Calibration moves the q90 coverage of the 20,641 matched rows from 0.798--0.828 to 0.879--0.887 and lowers the q90 pinball loss from 0.119--0.131 to 0.109--0.114; coverage stays below the nominal 0.90 and says nothing about censored days.

\subsection{Mechanism transfer and robustness}
\label{sec:transfer}

Table~\ref{tab:transfer} fixes the pooled-conformal policy at $\beta=0.75$ and varies the history used for training and calibration. The transfer gap $\Delta_{m\to h}=\overline{C}_{m\to h}(\mathcal I)-\overline{C}_{m\to m}(\mathcal I)$ compares the mean cost on an unseen evaluation history $h$ with the matched cost on the same rows $\mathcal I$. The largest gap arises for models trained on A and evaluated on C: 0.7303 against 0.6831, a gap of $+0.0472$. The C-trained models cost less on the A and B histories than on their own, the hardest one, so a negative gap signals an easier history rather than better transfer. All intervals are 95\% percentile intervals from 200 cluster-bootstrap replicates~\cite{efron1979bootstrap} over the 866 city--store clusters, with the same draws for both sides of a contrast; dates are not resampled separately and no multiplicity adjustment is applied, so we read the intervals as descriptive.

The reversal holds in all nine train/test cells. When excess stock is more costly, the raw-sales LightGBM with residual quantiles is the cheapest full-sample policy in every cell, with pooled conformal 0.0007 to 0.0323 more expensive; when shortage is more costly, pooled conformal is the cheapest in every cell, 0.0441 to 0.1298 below the raw-sales policy; at equal costs four cells go to each and one to the uncalibrated demand quantiles. Over the three settings, the largest transfer gap of pooled conformal, $+0.0472$, is the smallest worst case among the six policies whose forecasts read the evaluation history; the raw-sales quantile LightGBM reaches $+0.2258$, and Kaplan--Meier orders do not read that history. Over the seven final dates, with 2,507 to 3,385 donor rows each, the C history is the most expensive on every date, and the daily q90 coverage of pooled conformal ranges from 0.839 to 0.907. Under-coverage concentrates in series with at most ten complete fit-period days (0.752 to 0.769 over 532 rows) and on promotion days (0.854 to 0.865). These breakdowns are descriptive, not pre-specified.

\subsection{Real recorded histories}
\label{sec:real}

Applied to the same 35,000 real final-week rows, including rows with flagged stockouts, the q90 forecasts of the A-, B-, and C-trained pooled-conformal models give definite-miss lower bounds~\eqref{eq:miss} of 0.0895 $[0.0852, 0.0942]$, 0.0854 $[0.0808, 0.0896]$, and 0.0724 $[0.0684, 0.0766]$. A smaller bound can come from a higher forecast, so it does not by itself indicate better calibration. The bound is highest for rows with 1--4 recorded stockout hours (0.1465 to 0.1797) and lowest for 9--16 hours (0.0298 to 0.0334): when a product sells out early, recorded sales seldom exceed the forecast even when demand does. Only the complete-day stratum has an exact proxy.
